\section{XGBoost Tuning}
\label{sec:xgboost_tuning_results}

The retained market-only XGBoost model is refined through the controlled tuning procedure described in Section~\ref{sec:nonlinear_models}. Three configurations are compared: the initial model with a learning rate of 0.05 and depth three, a faster-learning specification with a rate of 0.10 and the same depth, and a deeper specification with a rate of 0.10 and depth four. Table~\ref{tab:xgboost-tuning-comparison} reports their training and validation performance. The selected configuration is shown in bold.

\begin{table}[!htbp]
\centering
\footnotesize
\setlength{\tabcolsep}{3.5pt}
\caption{Validation performance of the controlled XGBoost configurations.}
\label{tab:xgboost-tuning-comparison}
\begin{tabular}{rrrrrrrr}
\toprule
$\eta$ & Depth & Best round & $\mathrm{MAE}_{\mathrm{train}}$ & $\mathrm{MAE}_{\mathrm{val}}$ & $\mathrm{RMSE}$ & $R^2$ & $\overline{\rho}$ \\
\midrule
0.05 & 3 & 176 & 0.0836 & 0.0950 & 0.1317 & 0.466 & 0.749 \\
\textbf{0.10} & \textbf{3} & \textbf{90} & \textbf{0.0836} & \textbf{0.0951} & \textbf{0.1317} & \textbf{0.466} & \textbf{0.748} \\
0.10 & 4 & 54 & 0.0835 & 0.0952 & 0.1316 & 0.466 & 0.748 \\
\bottomrule
\end{tabular}
\end{table}


\noindent
The three configurations achieve nearly identical validation performance. Figure~\ref{fig:xgboost-base-selected-learning-curve} compares the initial and selected learning curves. Increasing the learning rate from 0.05 to 0.10 reduces the validation-selected number of boosting rounds from 176 to 90 without materially changing validation error, while increasing tree depth provides no further gain. The depth-three model with a learning rate of 0.10 is therefore retained for its faster convergence and lower complexity. Its complete parameter specification is reported in Appendix~\ref{app:xgboost_tuning_details}.

\begin{figure}[!htbp]
    \centering
    \includegraphics[width=0.82\textwidth]{04_02_xgboost_base_vs_selected_learning_curve.pdf}
    \caption[Learning curves for the initial and selected XGBoost configurations]{Training and validation MAE for the initial and selected market-only XGBoost configurations. The initial curves are semi-transparent, and the vertical lines mark each configuration's validation-selected boosting round.}
    \label{fig:xgboost-base-selected-learning-curve}
\end{figure}
